\documentclass[a4paper, 12pt]{report}
\usepackage[bottom=3cm]{geometry}
\usepackage{polski}
\usepackage{icomma}
\usepackage[utf8]{inputenc}
\usepackage[T1]{fontenc}
\usepackage{amsmath}
\usepackage{graphicx}
\graphicspath{ {graphics/} }
\usepackage{pgfplots}
\pgfplotsset{compat=newest}
\usepgfplotslibrary{groupplots}
\usepgfplotslibrary{dateplot}
\pgfplotsset{every axis/.append style={
             tick label style={/pgf/number format/fixed},
             font=\scriptsize,
             ylabel near ticks,xlabel near ticks,
             grid=major}}
\usepackage{diagbox}
\usepackage[justification=centering]{caption}
\usepackage{subcaption}
\usepackage{setspace}
\usepackage{array}
\usepackage{fancyhdr}
\pagestyle{fancy}
\usepackage[hidelinks]{hyperref} 
\usepackage[shortlabels]{enumitem}
\usepackage{siunitx}
\sisetup{locale = FR}

\usepackage{minted}
\renewcommand\listoflistingscaption{Spis listingów}
\usepackage[polish]{babel}
\usepackage{booktabs}
\usepackage{multirow}
\usepackage{rotating}
\usepackage{texlogos}
\pgfplotsset{compat=1.8}
\usepackage[toc,page]{appendix}
\renewcommand{\appendixtocname}{Dodatki}
\renewcommand{\appendixpagename}{Dodatki}

\newcolumntype{C}[1]{>{\centering\arraybackslash}p{#1}}

\fancyhead{}
\fancyhead[RO,LE]{\chaptername: \rightmark}

% Forbid to split given words
\hyphenation{LabVIEW}

\pgfplotsset{grid style={dashed}}
% This is necessary to plot softmax function
\pgfmathdeclarefunction{sumexp}{3}{%
\begingroup%
\pgfkeys{/pgf/fpu,/pgf/fpu/output format=fixed}%
\pgfmathsetmacro{\myx}{#1}%
\pgfmathtruncatemacro{\myxmin}{#2}%
\pgfmathtruncatemacro{\myxmax}{#3}%
\pgfmathsetmacro{\mysum}{0}%
\pgfplotsforeachungrouped\XX in {\myxmin,...,\myxmax}%
{\pgfmathsetmacro{\mysum}{\mysum+exp(\XX)}}%
\pgfmathparse{\mysum+exp(#1)}%
\pgfmathsmuggle\pgfmathresult\endgroup%
}%

\begin{document}

\newpage
\thispagestyle{empty}
\newgeometry{top=2.5cm, bottom=2.5cm, left=3cm, right=2.5cm}
\begin{onehalfspacing}
\begin{center}
	% TODO: vspaces are coarsely approximated, correct them
	\includegraphics[scale=0.75]{polsl_logo_bw_pl}
	\vspace{0.8cm}
	
	\fontsize{18}{18} \selectfont
	\textbf{\textsc{Politechnika Śląska \linebreak
	Wydział Automatyki, Elektroniki i~Informatyki}}
	\vspace{1.3cm}
	
	\fontsize{16}{16} \selectfont
	Projekt inżynierski
	\vspace{1.7cm}
	
	\fontsize{14}{14} \selectfont
	Wykorzystanie sieci neuronowych do detekcji ziaren
	w~obrazach termowizyjnych
	\vspace{5cm}
	
	\begin{flushleft}
	Autor: Maciej Ziaja \linebreak
	Kierujący pracą: dr inż. Sebastian Budzan \linebreak
	\end{flushleft}
	
	\vfill
	\fontsize{12}{12} \selectfont
	Gliwice, styczeń 2020
\end{center}
\end{onehalfspacing}
\restoregeometry


\pagenumbering{gobble}

% Introductory pages
\begin{abstract}
Tematem pracy jest klasyfikacja (rozpoznawanie typu) ziaren rud miedzi na
podstawie nagrań metodą termowizji aktywnej.
Zaproponowano sposób wykonania pomiarów oraz zbierania danych z~procesu
stygnięcia badanych próbek miedzi.
W~celu budowy zbioru danych przeprowadzono eksperymenty z~użyciem
kamery termowizyjnej wyposażonej w~makroobiektyw.
Zgromadzone materiały przeanalizowano pod kątem użyteczności w~zadaniu
klasyfikacji.
Zaproponowano możliwe sposoby ekstrakcji cech charakterystycznych
rud miedzi oraz wybrano najlepszy, polegający na detekcji detali obrazu
o~małej emisyjności cieplnej.
Przedstawiono metody zliczania i~śledzenia elementów charakterystycznych
rud miedzi, w~cyklu zdjęć.
Rozpatrywane sposoby ekstrakcji danych z~nagrań zaimplementowano w~języku
programowania Python.
Zgromadzone dane i~opracowane metody eksploracji danych wykorzystano
do konstrukcji prototypu sieci neuronowej klasyfikującej ziarna rud miedzi.
Przeprowadzono ocenę działania klasyfikatora oraz wytypowano najlepszy
uzyskany sposób na detekcję typu ziaren.
Walidacja uzyskanego modelu wskazała na jego użyteczność w~klasyfikacji
nagrań procesu stygnięcia rud miedzi.
Zaproponowano pomysły na rozbudowę stanowiska pomiarowego, automatyzację
procesu zbierania danych, perspektywy rozwoju projektu oraz rozpatrzono
możliwe alternatywne rozwiązania.
Stworzone oprogramowanie będzie stanowić ułatwienie i~bazę dla dalszych
badań ziaren rud miedzi metodą termowizyjną.
Efekty pracy mogą stanowić punkt wyjścia dla kolejnych projektów oraz prób
wdrożenia opracowanej technologii w~przemyśle.

\end{abstract}
 
\tableofcontents

\newpage
\pagenumbering{arabic}

% Chapters
\chapter{Wstęp}
Lorem ipsum dolor sit amet, consectetur adipiscing elit. Mauris condimentum odio non risus semper imperdiet. Duis eget ipsum fringilla, fermentum elit vitae, malesuada velit. Phasellus porttitor bibendum auctor. Aliquam non risus eget sapien viverra mollis ac non felis. In mattis egestas orci at pretium. In a felis vitae erat sollicitudin efficitur ac vel orci. Etiam aliquam placerat massa, sed varius nulla. Maecenas pulvinar, est eu condimentum facilisis, velit nisi viverra leo, eget gravida ipsum diam a purus. Nam feugiat leo purus, in aliquam quam eleifend sed. Etiam sed ligula sed neque consequat luctus vel eu ex.

Nam sagittis venenatis sagittis. Ut ac massa vehicula metus tempus sollicitudin. Donec aliquet eros vitae euismod dictum. Nam blandit iaculis felis vitae facilisis. In hac habitasse platea dictumst. Nam sagittis quam augue, at pretium lacus cursus at. Curabitur volutpat vulputate dui viverra pretium. Morbi vitae malesuada dolor. Cras hendrerit mi at interdum varius. Donec efficitur arcu sollicitudin urna luctus commodo sit amet quis odio.

Nullam auctor finibus enim, fermentum dictum velit. Nullam dictum, sapien id ullamcorper blandit, lorem risus egestas nisi, quis aliquam erat arcu ut risus. Ut ac tempus nunc, sit amet ornare erat. Sed sed nisi massa. Vivamus facilisis mi eros, eu malesuada libero placerat eget. Cras luctus aliquet nibh a pulvinar. Vestibulum id ex sed arcu ultricies accumsan in sit amet diam. Nam ut lacus vitae eros iaculis gravida. Nullam lobortis ante in tincidunt fermentum.

Aliquam et gravida nisi. Quisque risus felis, facilisis ullamcorper odio eget, rutrum mattis tellus. Praesent sit amet nibh elementum, rhoncus velit vitae, venenatis lectus. Sed pellentesque id nulla ut lobortis. Quisque et magna posuere, ornare ipsum a, molestie libero. Pellentesque a justo eget mauris cursus eleifend. Sed ac dolor et leo maximus malesuada. Nulla facilisi. Suspendisse sit amet molestie neque. Integer consectetur, neque id suscipit congue, orci tellus ultrices odio, eu lobortis turpis sapien ac eros. Sed porttitor lacus accumsan hendrerit facilisis. Morbi cursus nisl id auctor pretium.

Fusce sit amet dui eget lacus varius aliquam vitae vitae quam. Curabitur luctus est in elit convallis euismod. Ut id volutpat nibh. Fusce scelerisque luctus justo egestas maximus. Nam mollis libero id nulla porta pretium. Aliquam nisi odio, facilisis id nisl eget, dapibus vehicula odio. Integer tortor mi, tincidunt non dignissim vitae, laoreet ut erat. Curabitur velit purus, scelerisque non orci et, efficitur tempor ipsum. Donec dictum vitae nulla vel tincidunt. Aenean odio purus, egestas in aliquet sit amet, tempor sit amet metus. Sed vel tristique orci. Donec fringilla metus eget risus ultrices rhoncus. Donec gravida dui a scelerisque dictum.

\chapter{Założenia projektowe i~wykorzystane narzędzia}
\section{Analizowane ziarna rud miedzi} \label{sec:grains}
Zestaw analizowanych ziaren pochodzi z~projektu \textsc{sysmel} (\emph{System
sterowania procesem mielenia w~układzie z~młynem elektromagnetycznym})%
\footnote{Pełen tytuł projektu naukowego: \emph{Układ mielenia surowców
    mineralnych w~młynie elektromagnetycznym wraz z~systemem sterowania jego
    pracą, zapewniający wysoką efektywność technologiczną i~niską
    energochłonność w~zastosowaniach mikro i~makro-przemysłowych.}},
którego współtwórcami są naukowcy z~Politechniki Śląskiej.
Uzyskane próbki powstały w~procesie mielenia rud miedzi w~młynie
elektromagnetycznym.
Poszczególne klasy ziaren różnią się granulacją i~parametrami procesu ich
mielenia.
Spośród próbek, pochodzących z~projektu \textsc{sysmel}, wybrano cztery typy
ziaren, które będą przedmiotem zadania klasyfikacji.
Zdjęcia rozpatrywanych materiałów przedstawiono na rysunku \ref{fig:grains}.
\begin{figure}[htb]
    \hspace*{\fill}
    \begin{subfigure}[t]{0.35\textwidth}
        \centering
        \includegraphics[width=\textwidth]{E5R}
        \caption{Ziarna rudy miedzi klasy E5R}
    \end{subfigure}
    \hfill
    \begin{subfigure}[t]{0.35\textwidth}
        \centering
        \includegraphics[width=\textwidth]{E6R}
        \caption{Ziarna rudy miedzi klasy E6R}
    \end{subfigure}
    \hspace*{\fill}
    \vskip\baselineskip
    \hspace*{\fill}
    \begin{subfigure}[t]{0.35\textwidth}
        \centering
        \includegraphics[width=\textwidth]{E11R}
        \caption{Ziarna rudy miedzi klasy E11R}
    \end{subfigure}
    \hfill
    \begin{subfigure}[t]{0.35\textwidth}
        \centering
        \includegraphics[width=\textwidth]{E16R}
        \caption{Ziarna rudy miedzi klasy E16R}
    \end{subfigure}
    \hspace*{\fill}
    \caption{Próbki rozpatrywanych ziaren rud miedzi}
    \label{fig:grains}
\end{figure}

\section{Metody pomiarów termowizyjnych}
\label{sec:thermovision}
Pomiary termowizyjne (nazywane także \emph{termografią}) polegają na rejestracji 
promieniowania cieplnego obiektów w~celu ustalenia ich temperatury.
Badania z~wykorzystaniem termowizji można podzielić na dwie grupy:
\emph{termowizję pasywną} oraz \emph{termowizję aktywną}.

Technika pasywna polega na rejestracji obrazów obiektów bez ingerencji w~ich
temperaturę w~czasie trwania pomiarów.
Można w~ten sposób obserwować przepływ ciepła w~urządzeniach technicznych,
procesach przemysłowych oraz biologicznych.
W~rozważanym w~pracy przypadku technika ta nie ma jednak zastosowania.
W~warunkach pokojowych badany materiał ma w~całej swojej objętości podobną
temperaturę, co skutkuje obrazem równomiernego szumu w~rejestrowanym obrazie.

Bardziej zaawansowaną techniką są pomiary z~wykorzystaniem termowizji aktywnej.
W tym trybie pomiarów badane obiekty ogrzewa się w~powtarzalnych warunkach,
a~następnie rejestruje proces ich stygnięcia.
Podczas dostarczania ciepła do obiektu oraz jego oddawania na obrazach
termowizyjnych można zaobserwować strukturę badanego obiektu.
Przedmioty o~niejednorodnej i~złożonej budowie nagrzewają się oraz stygną
nierównomiernie, co jest rejestrowane przez kamery termowizyjne.
W~ten sposób można dostrzec cechy obiektu niewidoczne gołym okiem, takie jak 
zmiany jego gęstości, składu oraz uszkodzenia struktury.
Termowizję aktywną stosuje się szeroko w~badaniach naukowych oraz przemyśle.
Przykładowe aplikacje tej techniki to detekcja defektów w~produktach
przemysłowych oraz wykrywanie części konstrukcji podatnych na zużycie.
Termowizja aktywna ma charakter niedestruktywny oraz bezkontaktowy, co stanowi
jej zalety w~analizie materiałowej \cite{ciampa_thermography}.
Technika ta wymaga jednak etapowego procesu pomiarowego, potrzebne jest
przygotowanie instalacji grzewczej, a~nagrzewanie i~stygnięcie materiału może
być czasochłonne.
Przebieg zmian temperatury najlepiej rejestrować na materiałach wideo, aby
zmaksymalizować ilość danych zgromadzonych w~trakcie eksperymentu.

\section{Kamera termowizyjna \textsc{Flir~A320}}
\label{sec:camera}
Przedmiotem projektu jest analiza obrazów pochodzących z~przemysłowej kamery
termowizyjnej \textsc{Flir~A320}.
Firma \textsc{Flir} zajmuje się produkcją wysokiej jakości kamer i~detektorów do
celów profesjonalnych i~przemysłowych.
Używany aparat łączy wysoką jakość pomiaru z~nowoczesnymi funkcjami integracji
z~oprogramowaniem komputerowym.
Urządzenie komunikuje się z~komputerem za pomocą sieci internetowej, pozwalając
na kontrolę z~poziomu aplikacji oraz bibliotek programistycznych.
Kamera posiada także możliwość planowania automatycznych pomiarów i~alarmów,
oferuje również funkcje analityczne oraz wbudowany serwer internetowy
\cite{flir_a32x_manual}.

\subsection{Opis sprzętowy wykorzystywanej kamery}
Kamera ma postać podłużnego korpusu, do którego swobodnie można podłączać
obiektywy.
Na rysunku \ref{fig:camera}~przedstawiono zdjęcie wykorzystywanego urządzenia.
\begin{figure}[h]
    \centering
    \includegraphics[width=0.4\textwidth]{flir_a320}
    \caption{Kamera termowizyjna \textsc{Flir~A320} \cite{flir_camera_specs}}
    \label{fig:camera}
\end{figure}
Egzemplarz kamery znajdujący się w~laboratorium termowizji Politechniki Śląskiej 
został wyposażony w~obiektyw, pozwalający oglądać ziarna rud miedzi
w~powiększeniu.
Kamera cechuje się następującymi parametrami \cite{flir_camera_specs}:
\begin{itemize}
    \item typ detektora: niechłodzony mikrobolometr,
    \item rozdzielczość: 320 na 240 pikseli,
    \item częstotliwość odświeżania: \SI{9}{\hertz} do \SI{30}{\hertz},
    \item szerokość otworu: \textit{f}\num{1,3},
    \item autofokus z wbudowanym silnikiem,
    \item zakres pomiarowy temperatur:
          \begin{itemize}[label={$\diamond$}]
              \item od \SI{-15}{\celsius} do \SI{+50}{\celsius},
              \item od \SI{0}{\celsius} do \SI{350}{\celsius},
          \end{itemize}
    \item dokładność: \num{\pm2}\si{\celsius} lub \num{\pm2}\% odczytu,
    \item zakres wykrywanego widma promieniowania: \SI{7,5}{\micro\meter}
          do \SI{13}{\micro\meter}.
\end{itemize}

Kamera wykrywa temperaturę przez detektor zwany \emph{bolometrem}, który mierzy
energię niesioną przez fale elektromagnetyczne w~spektrum podczerwieni.
Kiedy fala pada na detektor aparatu, temperatura komórek matrycy rośnie
i~zwiększa się ich rezystancja elektryczna.
Wartości rezystancji są mierzone, a~na ich podstawie określana jest temperatura
\cite{vanhoof_infrared}.

Zakres temperatur kamery jest odpowiedni do przeprowadzenia eksperymentów
z~pomiarami ziaren rud miedzi metodą termowizji aktywnej.
Zaplanowano podgrzewanie próbek maksymalnie do temperatury około
\SI{80}{\celsius},~wartość ta mieści się w~zakresie pracy urządzenia.
Dokładność kamery jest zadowalająca, próbne materiały nagraniowe wskazały, że na
zdjęciach widocznych jest wiele szczegółów i~detali badanego materiału.
Przy klasyfikacji obrazów i~wzorców stygnięcia próbek jest to bardziej istotne
niż liczbowa dokładność pomiarowa przyrządu.

W~porównaniu ze zwykłymi, współczesnymi aparatami, rozdzielczość kamery
termowizyjnej może wydawać się mała.
Należy sobie jednak uzmysłowić, że w~standardowych aparatach piksele mają
rozkład Bayera, a~wartości składowych koloru są interpolowane.
W~kamerze termowizyjnej każda komórka matrycy dokonuje pełnego pomiaru wartości
temperatury.
Z~tego powodu bezpośrednie porównanie rozdzielczości używanego przyrządu
z~popularnymi aparatami może być mylące.
Oczywiście większa rozdzielczość kamery byłaby pożądana, jednak jej obecne
możliwości pozwalają na szczegółowe pomiary i~obserwację wielu detali ziaren
rud miedzi.

\subsection{Obiektyw kamery}
\label{subsec:lens}
Kamerę wyposażono w~obiektyw zbliżeniowy \textsc{Flir~T197415}, pozwalający
obserwować ziarna rud miedzi.
Jest to sprzęt zaprojektowany przez producentów używanej kamery i~przeznaczony
do pracy z~urządzeniami termowizyjnymi.
Przyrząd przedstawiono na rysunku \ref{fig:lens}.
\begin{figure}[h]
    \centering
    \includegraphics[width=0.5\textwidth]{flir_lens}
    \caption{Obiektyw \textsc{Flir~T197415} \cite{flir_lens_specs}}
    \label{fig:lens}
\end{figure}
Wybrany obiektyw jest przygotowany z~myślą o~obserwacji drobnych detali
powierzchni w~dużym zbliżeniu.
Używany model ma następujące parametry \cite{flir_lens_specs}:
\begin{itemize}
    \item ogniskowa: \SI{18,2}{\milli\meter},
    \item powiększenie: \num{1x1},
    \item pole widzenia: \SI{8}{\milli\meter} na \SI{6}{\milli\meter},
    \item odległość od płaszczyzny ostrzenia: \SI{20}{\milli\meter},
    \item głębia ostrości: \SI{0,3}{\milli\meter},
    \item przysłona: bez regulacji, równa otworowi systemu montażu kamery,
    \item budowa: trzy soczewki asferyczne.
\end{itemize}

Zgodnie ze specyfikacją producenta obiektyw ma powiększenie \num{1x1},~co może
wydawać się niedużą wartością.
Należy mieć jednak świadomość, że zwykłe obiektywy zmniejszają obraz padający na
matrycę.
Powszechnie jako granicę makrofotografii przyjmuje się powiększenie \num{1x1}.~%
W~obserwacji ziaren i~detali powierzchni nie jest jednak istotne powiększenie,
ale to że używany obiektyw jest \emph{zbliżeniowy}.
Oznacza to, że ma on bardzo małą odległość ostrzenia, czyli można go przysunąć
blisko obserwowanej powierzchni.
Obiektyw zbliżeniowy pozwala na obserwację z~dystansu
\SI{20}{\milli\meter},~typowe obiektywy ostrzą z~odległości parudziesięciu
centymetrów do ponad metra.
Dzięki temu obiektyw zbliżeniowy pozwala na obserwację bardzo drobnych detali
powierzchni.

\subsection{Oprogramowanie do obsługi kamery}
\label{subsec:camera_soft}
Jedną z~najważniejszych cech kamery jest łatwość integracji z~oprogramowaniem
komputerowym.
Kamerę można obsługiwać za pomocą programu \textsc{Flir} Tools%
\footnote{%
    Strona produktu programu \textsc{Flir} Tools:
    \url{https://www.flir.com/products/flir-tools/}
}.
Pozwala on na podgląd obrazu oraz wykonywanie zdjęć termowizyjnych.
Producent dostarcza również bibliotekę LabVIEW pozwalającą na zaawansowaną
pracę z~kamerą.
Do obsługi stanowiska został napisany program używający tych bibliotek, który
pozwala na nagrywanie materiałów wideo przy pomocy kamery.
Nagrania mają własnościowy format firmy \textsc{Flir}, można je odtwarzać
w~programie \textsc{Flir} Tools.
Program pozwala także na eksport stopklatek z~nagrania w~postaci plików
\textsc{jpeg}.
Przy obsłudze narzędzia ważne jest ustawianie zakresu temperatur na obrazie.
Wybrany zakres określa w~jaki sposób wartości temperatury są mapowane na kolory
na zdjęciu, zawarte w~tablicy \textsc{lut}.
Program oferuje tablice w~skali szarości, takie zostały użyte w~projekcie,
możliwy jest także wybór tablic w~postaci kolorowych gradientów.
Wybór zakresu wpływa na wygląd wyświetlanego obrazu oraz eksportowanych klatek.
Jego nieodpowiedni dobór może skutkować zbyt ciemnym, jasnym, lub mało
kontrastowym obrazem.
Aby zapewnić najlepsze wykorzystanie nagrań z~kamery oraz powtarzalny charakter
eksportu stopklatek korzystano z~opcji automatycznego doboru zakresu temperatur, 
jaki jest wbudowany w~program \textsc{Flir} Tools.

\section{Narzędzia programistyczne}

\subsection{Język programowania Python}
Założenia projektu wymagają użycia języka programowania pozwalającego na
zaawansowaną obróbkę obrazu oraz wydajne budowanie sieci neuronowych.
W~obu tych dziedzinach wiodącym językiem jest Python%
\footnote{%
    W~projekcie użyto języka Python 3,~który jest aktualnym standardem.
    Nie należy go mylić z~niekompatybilną i~przestarzałą wersją Python 2,~która
    traci wsparcie w~2020 roku.
}%
, posiadający bogaty zestaw bibliotek.
Język ten oferuje dużą wygodę programowania, co jest istotne przy pracy
badawczej.
Do wydajnych obliczeń numerycznych w~Pythonie służy biblioteka NumPy.
Oferuje ona klasy macierzy numerycznych oraz bogaty zestaw operacji
matematycznych.
Macierze biblioteki NumPy są mniej elastyczne od zwykłych list języka Python,
jednak oferują dużo większą wydajność obliczeniową, między innymi dzięki
wsparciu obliczeń wektorowych na odpowiednich procesorach.
Wiele innych bibliotek języka wykorzystuje pakiet NumPy, na przykład przy
przetwarzaniu obrazów, co pozwala na wysoką wydajność obliczeń.
Połączenie wygody programowania z~zadowalającą wydajnością bibliotek
numerycznych, sprawia że Python jest dobrym wyborem do realizacji założeń
projektu.

Język Python cechuje się także integracją z~notatnikami \emph{Jupyter}%
\footnote{Strona projektu Jupyter: \url{https://jupyter.org}.
}.
Notatniki stanowią interaktywny edytor kodu.
Ich zaletą jest możliwość programowania bezpośrednio wewnątrz dokumentu
formatowanego za pomocą języka \emph{Markdown}.
Taki sposób organizacji kodu ułatwia dokumentację i~prezentację wyników,
szczególnie w~przypadku prac o~charakterze badawczym i~naukowym.
Programy demonstrujące działanie algorytmów oraz przedstawiające analizę danych
opracowano w~formacie notatników Jupyter.

\subsection{Biblioteka przetwarzania obrazu Scikit-image}
Język Python oferuje bogaty zestaw bibliotek przetwarzania obrazu.
Najbardziej popularną z~nich jest biblioteka OpenCV napisana w~języku
C\texttt{++}.
Zdecydowano się jednak na wybór mniej znanego pakietu Scikit-image\footnote{%
\label{foot:version_req}Wersje używanych bibliotek są zawarte w~pliku
\texttt{requirements.txt} dołączonym do projektu.
Na jego podstawie można utworzyć wirtualne środowisko języka Python
z~odpowiednim zestawem bibliotek.
Więcej informacji o~strukturze projektu zawarto w~dodatku
\ref{ch:project_structure}.%
}.
Jest to biblioteka zorientowana na obliczenia naukowe i~badawcze oraz napisana
bezpośrednio z~myślą o~języku Python \cite{scikit-image}.
Scikit-image, w~porównaniu z OpenCV, oferuje bardziej spójny interfejs
programisty oraz lepiej wykorzystuje charakterystykę języka Python.
Dodatkowo wybrana biblioteka jest częścią zestawu Scikit, w~skład którego
wchodzi pakiet Scikit-learn służący do uczenia maszynowego.
Może on być przydatny przy tworzeniu sieci neuronowej, a~spójność bibliotek
z~pakietu Scikit jest niewątpliwie zaletą.
Wybrana biblioteka charakteryzuje się również dobrą dokumentacją
\cite{scikit_reference} oraz zestawem przykładów.

\subsection{Interfejs sieci neuronowych Keras}
\label{subsec:software_network}
Wybór bibliotek głębokiego uczenia w~języku Python jest bardzo szeroki.
Do najpopularniejszych pakietów należą: Keras, TensorFlow oraz PyTorch.
Zdecydowano się na wybór interfejsu biblioteki Keras%
\footnoteref{foot:version_req}.
Jest to pakiet nastawiony na elastyczność i~możliwość eksperymentowania
\cite{chollet_keras}.
Moduł Keras oferuje także rozbudowaną dokumentację oraz zestaw poradników
\cite{keras_docs}.
Keras nie jest jednak samodzielną biblioteką sieci neuronowych,
a~wyłącznie abstrakcyjnym interfejsem programistycznym.
Używanie go wymaga wyboru wewnętrznego silnika biblioteki.
Zdecydowano się na domyślną opcję użycia zaplecza pakietu TensorFlow.
Użyto wersji interfejsu Keras wchodzącej bezpośrednio w~skład modułu TensorFlow.
Taki sposób użycia bibliotek zapewnia ich najlepszą współpracę.



\chapter{Analiza i~ekstrakcja danych z~kamery termowizyjnej}
\section{Proces pomiarowy i~budowa zbioru danych}

\section{Analiza zebranych obrazów termowizyjnych}

\subsection{Prezentacja przykładowej serii pomiarowej}

\subsection{Przetwarzanie danych wizyjnych}

\subsection{Poprawa jakości obrazu}

\subsection{Automatyczny odczyt zakresu pomiarowego temperatur z~obrazu}
Jak wspomniano w~sekcji \ref{subsec:camsoft} jednym z~kluczowych czynników
decydujących o~wyglądzie obrazów pochodzących z~kamery jest zakres temperatur
mapowany na kolory w~obrazie.
Niestety aplikacja FLIR Tools nie pozwala na eksport zakresu temperatur wraz
ze zdjęciami w~formie liczbowej.
W~czasie zapisu zdjęć oprogramowanie dodaje na nich interfejs z~aktywną skalą
pomiarową, jednak jest on graficznie naniesiony na obraz.
Aby ułatwić w przyszłości pracę z~materiałami z~kamery opracowano dodatkowo
mechanizm ekstrakcji zakresu temperatur z~zdjęć pochodzących z kamery.

W~celu konstrukcji funkcji odczytu wartości liczbowych z~obrazu posłużono
się gotową siecią neuronową zaprojektowaną do detekcji tekstu na zdjęciach.
Zdecydowano się na użycie popularnej biblioteki \emph{Pytesseract}.
Aby poprawnie odczytać wartości z~obrazu najpierw przycięto je tak by w~kadrze
znajdowała się tylko odczytywana liczba.
Ponieważ przy eksporcie zdjęć program FLIR Tools nakłada interfejs na zdjęcia
w~identyczny sposób, kadrowanie obrazu jest takie same dla każdej próbki
pomiarowej.
Wycięte kadry są bardzo małej rozdzielczości, aby ułatwić sieci rozpoznawanie
liczb zdecydowano się przeskalować je w~górę.
W~czasie skalowania włączono mechanizm anty aliasingu aby wyrównać krawędzie
cyfr.
Ponieważ używana sieć uznaje za tło kolor biały oraz poszukuje liczb w~kolorze
czarnym barwy na zdjęciu odwrócono.
Następnie obraz poddano binaryzacji metodą \emph{otsu}.
Jest to popularna i~wydajna metoda binaryzacji, jej efektywność jest
maksymalna kiedy ilość pikseli tła oraz pierwszego planu jest zbliżona,
dlatego poprawne kadrowanie liczb sprzyja jakości ich binaryzacji\cite{sezgin}.
Na rysunku \ref{fig:temp_bounds} przedstawiono kolejne etapy przygotowania
obrazu do rozpoznania liczb.
Implementację opisanego mechanizmu odczytywania zakresu temperatur ze zdjęć
przedstawiono na listingu \ref{lst:temp_bounds}.

\begin{figure}[h]
	\centering
	\begin{subfigure}[h]{0.45\textwidth}
		\centering
		\includegraphics[width=\textwidth]{temp_bounds_scale}
		\caption{Przeskalowany kadr z~liczbą}
		\label{fig:temp_bounds_scale}
	\end{subfigure}
	\quad
	\begin{subfigure}[h]{0.45\textwidth}
		\centering
		\includegraphics[width=\textwidth]{temp_bounds_bin}
		\caption{Kadr z~liczbą po binaryzacji}
		\label{fig:temp_bounds_bin}
	\end{subfigure}
	
	\caption{Przygotowanie zakresu temperatur do odczytu przez sieć neuronową}
	\label{fig:temp_bounds}
\end{figure}

\begin{listing}[h]
\begin{minted}{python}
def get_temperature_bounds(img, bounds=(((6, 24), (283, 318)),
                                       ((219, 236), (283, 318)))):
    '''Extract temperature values from FLIR UI on image.'''
    img = invert(img)
    temp_txt = []
    for bound in bounds:
        bound_img = img[slice(*bound[0]), slice(*bound[1])]
        bound_img = rescale(bound_img, 4, anti_aliasing=True)
        thr = threshold_otsu(bound_img)
        img_txt = bound_img > thr
        img_txt = Image.fromarray(img_txt)
        temp = pytesseract.image_to_string(img_txt, config='digits')
        if temp is not '': 
            temp = float(temp)
        else:
            temp = 0
        temp_txt.append(temp)
    return temp_txt
\end{minted}
\caption{Funkcja języka Python do odczytywania zakresu temperatur ze zdjęć
         z~kamery}
\label{lst:temp_bounds}
\end{listing}


\section{Poszukiwanie zależności użytecznych w~klasyfikacji}

\subsection{Rozważane możliwości użycia cech obrazu w~klasyfikacji}

\subsection{Wybór algorytmu detekcji ziaren}

\subsection{Algorytm śledzenia ziaren w~serii zdjęć}

\subsection{Wizualizacja zebranych cech i~ocena ich użyteczności}


\chapter{Prototypowanie sieci neuronowej klasyfikującej ziarna miedzi}
\section{Konstrukcja prototypu sieci neuronowej}
Po zakończeniu procesu eksploracji danych można przystąpić do
konstrukcji sieci neuronowej odpowiedzialnej za rozpoznawanie ziaren rud
miedzi.
Przedstawiony problem wymaga klasyfikacji wieloklasowej wielowymiarowych
danych.
Problemy tego typu można rozwiązywać zarówno za pomocą klasycznych algorytmów
uczenia maszynowego jak i~uczenia głębokiego.
Zdecydowano się na wybór sieci neuronowych do klasyfikacji ziaren.
Jest to rozwiązanie najbardziej nowoczesne i~elastyczne.
Zgodnie z~opisem narzędzi programistycznych w~podsekcji
\ref{subsec:software_network}, proces budowania sieci oparto na interfejsie
Keras, z~zapleczem TensorFlow.

\subsection{Dobór struktury sieci} \label{subsec:nnbuild}
Kluczowym czynnikiem decydującym o~architekturze sieci jest posiadany
zbiór danych.
Wejściowy zestaw cech klas jest pięciowymiarowy, do klasyfikacji takich
danych wystarczająca powinna być standardowa sieć neuronowa.
Nie ma potrzeby używania bardziej złożonych sieci rekurencyjnych, czy
konwolucyjnych.

Projektowanie sieci neuronowych wymaga podjęcia szeregu wyborów, wśród nich
można wyróżnić następujące decyzje dotyczące:
\begin{itemize}
	\item ilości warstw sieci,
	\item typu warstw w sieci,
	\item ilości neuronów w~poszczególnych warstwach sieci,
	\item rodzajów funkcji aktywacji w~neuronach poszczególnych warstw sieci.
\end{itemize}
Definiując model należy rozpatrzeć przedstawione cechy budowanej sieci.
Nie istnieje jednoznaczny sposób bezpośredniego określenia najlepszego
klasyfikatora dla danego problemu.
Przy projektowaniu sieci należy kierować się przesłankami teoretycznymi,
doświadczeniem oraz testując i~porównując różne rozwiązania.
Przedstawiona konstrukcja oraz uzasadnienie jej wyboru bazuje na wynikach
walidacji, której mechanizm przedstawiono w~sekcji \ref{sec:validation}.
Szczegółowe dane, dotyczące porównań różnych konstrukcji i~parametrów
sieci przedstawiono w~dodatku \ref{ch:nn_comparison_table}.

Konstrukcję klasyfikatora rozpoczęto od wyboru ilości warstw.
W~standardowej konstrukcji każda sieć posiada warstwę wejściową oraz
wyjściową, których parametry są związane odpowiednio z~wektorem cech
oraz reprezentacją klas na wyjściu sieci.
Ponadto do aproksymacji dowolnej funkcji nieliniowej wystarczająca jest sieć
składająca się z~dwóch warstw.
Nie oznacza to jednak, że w~praktyce dwie warstwy wystarczą by rozwiązać
dowolny problem.
Aby zwiększać efektywność sieci pomiędzy warstwą wejściową oraz wyjściową
umieszcza się zestaw warstw nazywanych \emph{ukrytymi}.
Zazwyczaj posiadają one najwięcej neuronów spośród warstw w~sieci.
Mimo że jedna warstwa ukryta pozwala teoretycznie na rozwiązanie dowolnego
problemu, w~praktyce opłacalne jest stosowanie większej ilości warstw,
co poparte jest licznymi przykładami \cite{reed}.
Biorąc pod uwagę złożoność rozpatrywanego problemu zdecydowano się na
strukturę czterowarstwową.
Jest to rozmiar często spotykany w~sieciach tego typu, późniejsze testy
pokazały, że wybrana liczba warstw daje dobre rezultaty.
Sprawdzono, że zastosowanie jednej warstwy ukrytej powodowało pogorszenie
działania sieci.

Kolejnym krokiem projektowania sieci jest decyzja o~ilości neuronów
w~poszczególnych warstwach.
W przypadku pierwszej i~ostatniej warstwy jest to wartość prosta do
określenia.
Zaleca się, aby wejście sieci miało liczbę neuronów równą liczbie elementów
wektora cech, który w~rozpatrywanym przypadku ma długość równą pięć.
% TODO: cite
Ostatnia warstwa powinna mieć tyle neuronów ile wynosi liczba rozpoznawanych
klas, tak by każde wyjście sieci oznaczało prawdopodobieństwo przynależności
próbki do odpowiedniej klasy.
W~analizowanym zbiorze ziaren znajdują się cztery typy rud miedzi i~tyle
neuronów powinno znajdować się w~warstwie wyjściowej sieci.
W~warstwach ukrytych umieszczono większą liczbę neuronów, eksperymentalnie
stwierdzono, że sieć osiąga dobre wyniki dla 256 neuronów w~pierwszej
warstwie ukrytej i~128 w~drugiej.
Zmniejszenie liczby neuronów o~połowę w~kolejnej warstwie ukrytej jest
często spotykanym zabiegiem.
% TODO: cite

Następnie rozpatrzono dostępne funkcje aktywacji poszczególnych warstw.
Dla warstwy wyjściowej należy wybrać funkcję o~wartościach od zera do jeden.
Jest to uzasadnione faktem, że wyjścia powinny reprezentować
prawdopodobieństwo zaliczenia próbki do danej klasy.
Przy klasyfikacji zazwyczaj stosuje się funkcję softmax, jej wykres
przedstawiono na rysunku \ref{fig:softmax}.
\begin{figure}[htb]
\centering
\begin{tikzpicture}[trim axis left]
    \begin{axis}[title={$ Softmax(z) $}, width=8cm, height=6cm,
                 ylabel=$ \sigma(z) $, xlabel=$ z $, xmin=-5, xmax=5]
       \addplot[blue, domain=-5:5, samples=51] 
       {exp(x) / sumexp(x, -4, 0)};
    \end{axis}
\end{tikzpicture}
\caption{Funkcja aktywacji softmax}
\label{fig:softmax}
\end{figure}
W~przypadku pozostałych warstw wybór funkcji aktywacji jest mniej oczywisty.
Rozpatrzono trzy rodziny funkcji:
\begin{itemize}
	\item funkcję sigmoid,
	\item funkcje z~grupy jednostek liniowych (ang. \emph{linear unit}):
		\begin{itemize}	
        	\item ReLU (ang. \emph{Rectified Linear Unit}),
        	\item ELU (ang. \emph{Exponential Linear Unit}),
        \end{itemize}
	\item funkcję tangens hiperboliczny.
\end{itemize}
Rozpatrywane funkcje aktywacji przedstawiono na rysunku \ref{fig:activation}.
\begin{figure}[htb]
\centering
\begin{subfigure}[t]{0.3\textwidth}
\centering
\begin{tikzpicture}[trim axis left]
    \begin{axis}[title={$ Sigmoid(z) $}, width=5cm, height=4cm,
                 ylabel=$ \sigma(z) $, xlabel=$ z $, xmin=-5, xmax=5]
       \addplot[blue, domain=-5:5, samples=51] 
       {1 / (1 + e^(-x)};
    \end{axis}
\end{tikzpicture}
\caption{Funkcja aktywacji tangens hiperboliczny}
\end{subfigure}
\hspace{0.25cm}
\begin{subfigure}[t]{0.3\textwidth}
\centering
\begin{tikzpicture}[trim axis left]
    \begin{axis}[title={$ LU(z) $}, width=5cm, height=4cm,
                 ylabel=$ \sigma(z) $, xlabel=$ z $, xmin=-3, xmax=3,
                 legend style={at={(0.02, 0.98)}, anchor=north west}]
       \addplot[red, domain=-3:3, samples=31] 
       {max(0, x)};
       \addplot[blue, domain=-3:3, samples=31] 
       {x < 0 ? 0.1*(e^x - 1) : x};
    \legend{ReLU, ELU}
    \end{axis}
\end{tikzpicture}
\caption{Rodzina funkcji aktywacji jednostek liniowych}
\end{subfigure}
\hspace{0.25cm}
\begin{subfigure}[t]{0.3\textwidth}
\centering
\begin{tikzpicture}[trim axis left]
    \begin{axis}[title={$ Tanh(z) $}, width=5cm, height=4cm,
                 ylabel=$ \sigma(z) $, xlabel=$ z $, xmin=-5, xmax=5]
       \addplot[blue, domain=-5:5, samples=51] 
       {tanh(x)};
    \end{axis}
\end{tikzpicture}
\caption{Funkcja aktywacji tangens hiperboliczny}
\end{subfigure}
\caption{Dostępne funkcje aktywacji rozpatrywane do użycia w~warstwach
         ukrytych}
\label{fig:activation}
\end{figure}
Funkcja sigmoid jest często stosowaną, klasyczną funkcją aktywacji.
Jej wadą jest jednak zanikanie wartości pochodnej.
Obecnie częściej stosowane są funkcję z~rodziny jednostek liniowych,
szczególnie z wyciekiem, czyli małymi wartościami dla ujemnych argumentów.
Przetestowano również funkcję tangens hiperboliczny, która ma kształt
podobny do funkcji sigmoid.
Mimo, że obecnie funkcje jednostek liniowych są najbardziej popularne
w~rozpatrywanym przypadku sieć osiągała najlepsze wyniki dla funkcji tangens.

Opisany model sieci neuronowej należy zdefiniować za pomocą interfejsu
Keras.
Do inicjalizacji sieci o~standardowej liniowej strukturze służy funkcja
\mintinline{python}{Sequential()}, która może przyjąć listę warstw w~sieci.
Funkcja \mintinline{python}{Dense()} tworzy warstwę łącząca każdy neuron
z~każdym wyjściem poprzedniej warstwy.
Jej parametry definiują liczbę neuronów w~warstwie oraz ich funkcję aktywacji.
Funkcję języka Python implementują opisywaną sieć, za pomocą przedstawionych
elementów biblioteki Keras, przedstawiono na listingu \ref{lst:model}.
\begin{listing}[htb]
\begin{minted}{python}
def default_grain_classifier_model():
    '''
    Get default uncompiled model for grain classifcation,
    based on 5 step cooling process using number of blobs.
    '''
    model = keras.Sequential([
        keras.layers.Dense(5, activation='tanh'),
        keras.layers.Dense(256, activation='tanh'),
        keras.layers.Dense(128, activation='tanh'),
        keras.layers.Dense(4, activation='softmax')
    ])
    return model
\end{minted}
\caption{Funkcja języka Python definiująca model sieci neuronowej}
\label{lst:model}
\end{listing}

\subsection{Trening sieci neuronowej}
\label{subsec:train}
Kolejnym krokiem w~budowie klasyfikatora ziaren jest trening sieci neuronowej.
Jest to proces. w~którym wagi neuronów są dopasowywane do zbioru uczącego.
Rozpatrywany problem jest typu uczenia nadzorowanego.
Przygotowane dane mają znane etykiety klas, to znaczy są odgórnie przypisane
do danego typu.
Narzucone etykiety nazywane są \emph{wiedzą ekspercką}, ich prawdziwość jest
pewna.
Etykiety przypisane przez klasyfikator będą oceniane przez porównanie
z~etykietami eksperckimi.
Trening polega na zbliżaniu wyników działania sieci do etykiet zbioru
uczącego.

Przed rozpoczęciem uczenia sieci należy podzielić dostępne dane na zbiór
treningowy oraz~testowy.
W~skład zbioru testowego nie powinny wchodzić dane uczące.
Ocena zdolności sieci do generalizacji problemu może bazować tylko
na danych, które nie brały udziału w~jej treningu.
Aby ułatwić podział Biblioteka Scikit-learn udostępnia funkcję
\mintinline{python}{train_test_split()}, która zwraca
podane zbiory, podzielone na części treningowe i~testowe.
Odpowiednie wykorzystanie funkcji wymaga użycia jej parametrów opcjonalnych.
Argument \mintinline{python}{stratify} sprawia, że podzielone zbiory
zawierają takie same proporcje klas.
Taka metoda podziału nazywana jest \emph{losowaniem warstowym}.
Uaktywnienie tej opcji jest szczególnie ważne, ze względu na mały rozmiar
posiadanego zbioru danych.
Gdyby dane dzielić w~pełni przypadkowo istniałoby ryzyko nadmiernej
reprezentacji klasy w~danej grupie oraz jej braku w~innej.
Kolejnym istotnym parametrem jest \mintinline{python}{test_size}, który
decyduje o~rozmiarze zbioru testowego.
Ponieważ w~zbiorze są trzy egzemplarze każdej klasy odpowiednie jest
wydzielenie do testów jednej trzeciej próbek.
Ostatni parametr \mintinline{python}{random_state}, to ziarno generatora
losowego podziału zbiorów.
Aby móc porównywać działanie sieci, pomiędzy wielokrotnymi uruchomieniami
programu, należy wyeliminować z~niego czynniki przypadkowe i~podać
funkcji stałe ziarno, co zapewni powtarzalny podział zbioru.
Oczywiście wydzielenie w~narzucony sposób zbioru testowego, szczególnie
w~przypadku tak małej ilości danych, nie daje obiektywnej oceny modelu.
Jest on jednak wystarczający do budowy i~testowania pierwszego prototypu
sieci.
W~sekcji \ref{sec:validation} przedstawiono konstrukcję i~wyniki działania
bardziej miarodajnego procesu walidacji.

Aby móc przystąpić do treningu należy określić parametry uczenia
sieci.
Konfiguracja tego procesu odbywa się przez wywołanie metody
kompilacji.
Metoda przyjmuje parametry algorytmu optymalizacji sieci, miary błędu oraz
metryki.
Argument \mintinline{python}{optimizer} przyjmuje  nazwę stosowanego
algorytmu uczenia sieci.
Rozważono dwa popularne metody treningu: sgd (ang. \emph{stochastic gradient
descent)} oraz adam (ang. \emph{adaptive moment estimation}).
Metoda sgd jest najbardziej popularnym i~podstawowym sposobem uczenia, jednak
algorytm adam, jest rozwiązaniem nowszym, polecanym w~problemach klasyfikacji.
%TODO: cite
Zgodnie z~tymi przesłankami optymalizator adam okazał się dawać najlepsze
rezultaty i~to on został wybrany do uczenia sieci.
Parametr \mintinline{python}{loss} przyjmuje nazwę sposobu liczenia błędu
w~sieci.
Przykładem miary błędu używanej w~sieciach neuronowych jest popularny
błąd średniokwadratowy.
Dla problemów klasyfikacji lepszy jest jednak błąd obliczany metodą
\emph{rzadkiej kategoryzacyjnej entropii krzyżowej} (ang. \textit{sparse
categorical crossentropy}).
Jest to złożona metoda wykorzystująca funkcję logarytmiczną.
Metoda kategoryzacyjnej entropii okazała się najlepszą funkcją błędu
dla analizowanego przypadku.
Ostatni parametr metryki to wartość obliczana w~celu oszacowania poprawności
działania sieci.
Wybrano standardową metrykę dokładności sieci, czyli stosunku poprawnie
zakwalifikowanych próbek do wszystkich analizowanych.

Ostatnim etapem uczenia jest realizacja treningu sieci, dokonywana za
pomocą metody \mintinline{python}{model.fit()}.
Metoda przyjmuje argumenty zbioru treningowego wraz z~zestawem etykiet
oraz liczbę epok treningu, zwraca obiekt zawierający przebieg uczenia sieci.
Na rysunku \ref{fig:history} przedstawiono dokładność i~błąd w~kolejnych
epokach treningu.
Na ich podstawie określono, że odpowiednia liczba epok wynosi 300.
\begin{figure}[htb]
	\centering
	\begin{subfigure}[t]{0.3\textwidth}
		\centering
		\input{figures/neural_network_trainig_all}
		\caption{Przebieg treningu sieci dla metody zliczania wszystkich
		         ziaren}
	\end{subfigure}
	\hfill
	\begin{subfigure}[t]{0.3\textwidth}
		\centering
		\input{figures/neural_network_trainig_remaining}
		\caption{Przebieg treningu sieci dla metody zliczania śledzonych 
		         ziaren}
	\end{subfigure}
	\hfill
	\begin{subfigure}[t]{0.3\textwidth}
		\centering
		\input{figures/neural_network_trainig_percent}
		\caption{Przebieg treningu sieci dla metody procentowego zliczania
		         śledzonych ziaren}
	\end{subfigure}
	\caption{Przebieg treningu sieci dla różnych metod zliczania ziaren}
	\label{fig:history}
\end{figure}
Opisany plan treningu sieci realizuje kod przedstawiony na listingu
\ref{lst:train}.
\begin{listing}[htb]
\begin{minted}{python}
    X_train, X_test, y_train, y_test = train_test_split(X, y,
        stratify = y,
        test_size = 0.33,
        random_state = 1)

    model = default_grain_classifier_model()

    model.compile(optimizer='adam',
                  loss='sparse_categorical_crossentropy',
                  metrics=['accuracy'])
    
    history = model.fit(X_train, y_train, epochs=256)
\end{minted}
\caption{Kod treningu sieci neuronowej klasyfikującej ziarna miedzi}
\label{lst:train}
\end{listing}

Działanie sieci można sprawdzić na zbiorze testowym.
Należy mieć na uwadze, że przy małej ilości danych i~przedstawionym
podziale na zbiory treningowy oraz testowy nie będzie to test miarodajny.
Jest on jednak akceptowalny przy pierwszych testach sieci podczas jej
konstrukcji.
Wyniki działania sieci przedstawia tabela \ref{tab:blobtest}.
Przedstawione wartości błędu oraz dokładności potwierdzają przypuszczenia
na temat użyteczności różnych metod zliczania detali, które przedstawiono
w~podsekcji~\ref{subsec:datavis}.

\begin{table}[htb]
	\centering
	\begin{tabular}{c|C{2cm}|C{2cm}}
	\toprule
	\multirow{2}{*}{metoda zliczania detali} & \multicolumn{2}{c}{wskaźnik} \\ 
                                         & błąd       & dokładność      \\ \midrule
wszystkie detale              & 1,45       & 0,25            \\
śledzone detale              & 7,50       & 0,5             \\
procent śledzonych detali   & 0,63       & 0,75           \\   
	\bottomrule
	\end{tabular}
\caption{Wskaźniki oceny działania sieci na zbiorze testowym}
\label{tab:blobtest}
\end{table}

\section{Walidacja i~ocena działania sieci} \label{sec:validation}
Po stworzeniu prototypu klasyfkiatora przedstawionego w~rozdziale
\ref{subsec:nnbuild} należy przeprowadzić walidację zbudowanej sieci
neuronowej.
Najpopularniejszą metodą miarodajnej walidacji sieci jest 
\emph{k-krotny sprawdzian krzyżowy} (ang. \textit{k-fold cross-validation}).
Metoda ta polega na podziale dostępnych danych na k części, i~kolejnym
wydzielaniu jednej z~części danych jako zbioru testowego.
Walidacje powtarza się k-krotnie, tak by każda część była wykorzystana
jako dane testowe.
Na końcu testu wyniki są uśrednianie. 
Dzięki temu sprawdzian krzyżowy daje dobry pogląd na ogólną zdolność sieci do
generalizacji rozwiązywanego problemu.

Często po procesie walidacji dokonuje się testu sieci przy użyciu osobnego
zestawu danych.
Niestety posiadany zbiór jest zbyt mały, aby wydzielić z niego osobny zestaw
testowy.
Z~tego powodu zdecydowano się oceniać sieć jedynie na podstawie metody
sprawdzianu krzyżowego.

Biblioteka Keras nie posiada wbudowanych funkcji zaawansowanych technik
walidacji, dlatego należy przygotować je samodzielnie.
W~tym celu użyto biblioteki Scikit-learn, która oferuje funkcje podziału
zbiorów danych.
Jedną z~nich jest funkcja zwracająca iterowalny zestaw k-krotnych
podziałów zbioru danych.
Jak przedstawiono w~sekcji \ref{subsec:train}, ze względu na mały rozmiar
zbioru danych, przy podziale należy zastosować losowanie warstwowe.
Na listingu \ref{lst:crossval} przedstawiono funkcję pozwalającą na walidację
modelu klasyfikatora biblioteki Keras z~użyciem sprawdzianu krzyżowego.
Funkcja przyjmuje argumenty w~postaci skompilowanego modelu sieci oraz
zbioru danych i~etykiet.
Zwracana jest macierz, w~której rzędy oznaczają wyniki testów dla kolejnych
podziałów zbioru danych.
\begin{listing}[htb]
\begin{minted}{python}
def network_cross_validation(model, X, y, n_splits):
    '''Compute cross validation fold scores for given keras model.'''
    eval_scores = []
    
    folds = StratifiedKFold(n_splits=n_splits).split(X, y)
    for train_index, test_index in folds:
        x_train, x_test= X[train_index], X[test_index]
        y_train, y_test= y[train_index], y[test_index]
        
        model.fit(x_train, y_train, epochs=300, verbose=0)
        eval_scores.append(model.evaluate(x_test, y_test, verbose=0))
    return eval_scores
\end{minted}
\caption{Funkcja języka Python definiująca model sieci neuronowej}
\label{lst:crossval}
\end{listing}
Sposób wykorzystania zbudowanej funkcji walidacji modelu przedstawia
listing \ref{lst:val}.
Na końcu procesu oceny sieci wyniki kolejnych kroków sprawdzianu krzyżowego
należy uśrednić.
Efekty walidacji dla trzech sposobów zliczania detali na obrazach przedstawia
tabela \ref{tab:blobval}.

Zgodnie z~wcześniejszymi przewidywaniami metoda procentowego zliczania
detali na obrazach okazała się dawać najlepsze rezultaty.
Dokładność na poziomie 91\% jest wystarczająca do uznania sieci za dobry
klasyfikator.
Pozostałe metody zliczania okazały się mniej precyzyjne, co wskazuje że należy
je odrzucić.

Przedstawiony mechanizm walidacji wykorzystano do dostrojenia parametrów
sieci opisywanych w~podeskcji \ref{subsec:nnbuild}.
Uzasadnienie doboru przedstawionej w~toku pracy struktury sieci wynika z~prób
maksymalizacji dokładności oblicznej przy pomocy sprawdzianu krzyżowego.
W~czasie prototypowania sieci dokonano liczbych porównań, aby znaleźć
najlepszą strukturę.
Szczegółowe zestaweienie rozpatrywanych konfiguracji i~wartości parametrów
załączono w~dodatku \ref{ch:nn_comparison_table}.

\begin{listing}[htb]
\begin{minted}{python}
X = np.array(X)
y = np.array(y)

model = default_grain_classifier_model()
model.compile(optimizer='adam',
              loss='sparse_categorical_crossentropy',
              metrics=['accuracy'])

score = network_cross_validation(model, X, y, 3)

print('Folds scores: (loss, acc)\n', score)
score = np.array(score)
print('Cross validation mean score (loss, acc):\n',
      score.mean(axis=0), '\n')
\end{minted}
\caption{Wykorzystanie funkcji sprawdzianu krzyżowego do oceny działania
         sieci}
\label{lst:val}
\end{listing}

\begin{table}[htb]
	\centering
	\begin{tabular}{c|C{2cm}|C{2cm}}
	\toprule
	\multirow{2}{*}{metoda zliczania detali} & \multicolumn{2}{c}{wskaźnik} \\ 
                                         & błąd       & dokładność      \\ \midrule
wszystkie detale             & 6,59       & 0,41            \\
śledzone detale              & 3,99       & 0,50             \\
procent śledzonych detali   & 1,37       & 0,91          \\   
	\bottomrule
	\end{tabular}
\caption{Wskaźniki oceny działania sieci uzyskane metodą sprawdzianu
         krzyżowego}
\label{tab:blobval}
\end{table}

Dodatkowo opracowano kod generujący \emph{macierz pomyłek} (ang.
\textit{confusion matrix}).
Jest to tablica porównująca wiedzę eksperta z~wynikami klasyfikacji na zbiorze
testowym.
Wiersze macierzy oznaczają poprawną klasyfikację kolejnych próbek, zgodnie
z~etykietami eksperckimi.
Kolumny tabeli to typy klasyfikatora.
Przekątna tabeli zawiera informacje o~elementach poprawnie sklasyfikowanych,
etykiety klasyfikatora są zgodne z~danymi eksperckimi.
Komórki spoza przekątnej oznaczają błędną klasyfikację.
Obserwacja macierzy pozwala zaobserwować, które typy ziaren są ze sobą mylone
w~trakcie klasyfikacji.

Ze względu na małą ilość posiadanych danych macierz pomyłek opracowano
na podstawie mechanizmu sprawdzianu krzyżowego.
Na każdym jego etapie budowano macierz pomyłek, a~ostatecznie tablicę
uśredniono.
Wartości w~tabeli przedstawiono w~stosunku procentowym.
Funkcję generującą uśrednioną macierz pomyłek przedstawia listing
\ref{lst:mean_confusion_matrix}.
\begin{listing}[htb]
\begin{minted}{python}
def mean_confusion_matrix(model, X, y, n_splits):
    '''
    Compute mean confusion matrix using cross
    validation with n splits.
    '''
    conf_matrix = np.zeros((4, 4))

    folds = StratifiedKFold(n_splits = 3).split(X, y)
    for train_index, test_index in folds:
        x_train, x_test= X[train_index], X[test_index]
        y_train, y_test= y[train_index], y[test_index]
        
        model.fit(x_train, y_train, epochs=300, verbose=0)
        y_pred = model.predict_classes(x_test)

        for test, pred in zip(y_test, y_pred):
            conf_matrix[test][pred] = conf_matrix[test][pred] + 1
    
    return conf_matrix / n_splits
\end{minted}
\caption{Funkcja języka Python generująca uśrednioną macierz pomyłek}
\label{lst:mean_confusion_matrix}
\end{listing}

Zdecydowano się wygenerować macierz pomyłek dla najlepszego sposobu
klasyfikacji, bazującego na procentowym śledzeniu detali.
Wyniki zestawiono w~tabeli \ref{tab:confusion_matrix}.
Zgodnie z~wynikami sprawdzianu krzyżowego, większość elementów tablicy znajduje
się na jej przekątnej.
To znaczy, że wszystkie próbki klas E5R, E11R oraz E16R są poprawnie
klasyfikowane.
Widoczna jest pomyłka sieci polegająca na niewłaściwej klasyfikacji próbki
klasy E6R jako E16R.

\begin{table}[htb]
\centering
\begin{tabular}{cccccc}
                                 &                            & \multicolumn{4}{c}{Wiedza klasyfikatora} \\ \cmidrule[1pt]{2-6}
                                 & \multicolumn{1}{c|}{klasa} & E5R      & E6R      & E11R     & E16R    \\ \cmidrule{2-6}
\multirow{4}{*}{\rotatebox[origin=c]{90}{\parbox{2cm}
{\centering Wiedza eksperta}}} & \multicolumn{1}{c|}{E5R}   & 1,00     & 0,00     & 0,00     & 0,00    \\
                                 & \multicolumn{1}{c|}{E6R}   & 0,00     & 0,67     & 0,00     & 0,33    \\
                                 & \multicolumn{1}{c|}{E11R}  & 0,00     & 0,00     & 1,00     & 0,00    \\
                                 & \multicolumn{1}{c|}{E16R}  & 0,00     & 0,00     & 0,00     & 1,00 \\
                                 \cmidrule[1pt]{2-6}
\end{tabular}
\caption{Tabela pomyłek najlepszego klasyfikatora}
\label{tab:confusion_matrix}
\end{table}


\chapter{Podsumowanie} \label{ch:summary}
Wykonany projekt zrealizował cel jakim była konstrukcja klasyfikatora ziaren
rud miedzi, z~użyciem termowizji aktywnej oraz sieci neuronowych.
Zaproponowane sposoby zbierania danych, ekstrakcji cech oraz model sieci
okazały się wystarczające do osiągnięcia ponad dziewięćdziesięcio-procentowej
dokładności.
Jest to zadawalający wynik, świadczący o~dobrej jakości klasyfikacji.
W~toku projektu wyciągnięto liczne wnioski dotyczące alternatywnych
rozwiązań oraz możliwości rozwoju projektu.

Skonstruowany prototyp sieci dobrze pełni rolę klasyfikatora rud miedzi.
Zbudowane mechanizmy walidacji modelu mogą być pomocne przy jego
przyszłej rozbudowie oraz ocenie alternatywnych struktur sieci neuronowych.
Ze względu na bezkontaktowy i~niedestruktywny charakter pomiarów
termowizją aktywną, przedstawiono sposób klasyfikacji może znaleźć liczne
zastosowania.
Opracowane algorytmy mogą być podstawą do wykorzystania sieci neuronowych
do klasyfikacji ziaren rud miedzi w~środowisku przemysłowym.
Połączenie termowizji aktywnej oraz uczenia maszynowego może znaleźć
zastosowanie w~kontroli jakości półproduktów zawierających rudy miedzi.
Skonstruowane funkcje przetwarzania obrazów z~kamery, automatyczne kadrowanie
oraz odczytywanie temperatur ze zdjęć mogą być pomocne w~przyszłej pracy
z~wykorzystywanym sprzętem.
Zgromadzone materiały, techniki pomiarowe i~algorytmy mogą być pomocne
w~przyszłych pracach naukowych oraz dyplomowych bazujących na wykorzystaniu
sieci neuronowych w~klasyfikacji ziaren różnorakich materiałów.

Jak wspomniano w~sekcji \ref{sec:meas} proces budowy bazy danych był
czasochłonny oraz niepozbawiony niedokładności.
Zgromadzony zestaw nagrań okazał się wystarczający do spełnienia zadania
klasyfikacji rud miedzi, jednak możliwe są liczne usprawnienia stanowiska
pomiarowego.
Zgodnie z~opisem wykorzystywanej kamery przedstawionym w~sekcji
\ref{sec:camera}, istnieją liczne narzędzia pozwalające na automatyzację
procesu nagrywania materiałów wideo.
Wykorzystywany aparat został wyposażony w~bogaty zestaw sposobów komunikacji
z~komputerem oraz bibliotekę LabVIEW.
Istnieje zatem możliwość realizacji projektu rozbudowy stanowiska
laboratoryjnego termowizji w~celu automatyzacji procesu ogrzewania oraz
nagrywania próbek.
Pozwoliłoby to na znacznie szybsze utworzenie dużego zestawu danych oraz
poprawiłoby jakość zbieranych materiałów.
Budowa bardziej stabilnego i~zautomatyzowanego stanowiska pomiarowego
wyeliminowałaby ryzyko utraty ostrości obrazu pochodzącego z~kamery.
Dodatkowo rozbudowa stanowiska pozwoliłaby na realizację alternatywnej
metody pomiarów, polegającej na ogrzewaniu próbek do osiągnięcia zadanej
temperatury.

Zebrane materiały wideo stanowią zawierają wiele szczegółów oraz danych
charakteryzujących badane materiały.
W~podsekcji \ref{subsec:featureextr} rozpatrzono możliwe sposoby ekstrakcji
cech ze zgromadzonych nagrań.
Wybrana metoda śledzenia detali o~małej emisyjności okazała się odpowiednia
do wykorzystania jej w~klasyfikacji ziaren.
Przeprowadzona analiza świadczy, że uzyskany zestaw cech próbek dobrze
oddaje charakter fizykalnego procesu stygnięcia materiałów.
Opracowana technika ekstrakcji danych jest efektywna i~nie wymaga użycia
skomplikowanych modeli sieci neuronowych.
Rozpatrzono także alternatywne rozwiązania, które mogą również okazać się
skuteczne i~interesujące.
Po rozbudowie stanowiska oraz zebraniu większej ilości danych warta 
rozpatrzenia może okazać się możliwość wykorzystania nowoczesnych
konwolucyjnych sieci neuronowych.
Zebrane w~pracy materiały i~wnioski mogą posłużyć do realizacji projektów,
które rozwiną temat badania ziaren metodą termowizji.

Projekt spełnił swoje założenia, a~jego efektem jest działający klasyfikator
rud miedzi.
Wymienione wnioski mogą prowadzić do jego rozwoju oraz testów alternatywnych
rozwiązań.
Projekt może być podstawą do dalszych badań naukowych, jak i~prób wdrożenia
w~praktycznych zastosowaniach.


\begin{appendices}
\chapter{Porównanie konfiguracji i~parametrów modelu sieci}
\label{ch:nn_comparison_table}
\begin{table}[!h]
    \centering
    \begin{tabular}{c|c|c|c}
        \toprule
        Parametr & Wartość   & Błąd & Dokładność \\
        \midrule
        \multirow{4}{*}{Funkcja aktywacji}                 
            & sigmoid   & 1,28 & 0,58       \\
            & relu      & 1,71 & 0,83       \\
            & elu       & 2,04 & 0,92       \\
            & tanh      & 1,44 & 0,92       \\
        \midrule
        \multirow{2}{*}{Liczba warstw ukrytych}
            & 1         & 1,71 & 0,67       \\
            & 2         & 1,22 & 0,92       \\
        \midrule
        \multirow{3}{*}{\parbox{5cm}
            {\centering Liczba neuronów w~warstwach ukrytych}}
            & 128 i 64  & 1,47 & 0,83       \\
            & 256 i 128 & 1,42 & 0,92       \\
            & 512 i 256 & 0,97 & 0,92       \\
        \midrule
        \multirow{2}{*}{Algorytm uczenia}            
            & sgd       & 0,96 & 0,58       \\
            & adam      & 0,68 & 0,92       \\
        \bottomrule
    \end{tabular}
\end{table}

\chapter{Organizacja projektu}
\label{ch:project_structure}
\section{Struktura plików projektu}
Podczas pisania funkcji języka Python dbano o~ich zwięzłość, dobry podział
realizowanych zadań oraz ograniczenie długości bloków kodu.
Dzięki temu program zachowuje zasady dobrego programowania \cite{martin_code}.
Każdą funkcję opatrzono w~specjalny komentarz \emph{docstring} będący standardem
w~świecie języka Python.
Funkcje umieszczono w~odpowiednich plikach, które zawierają także przykłady ich
użycia i~obrazują działanie oraz porównują efekty.
Struktura plików w~projekcie jest następująca:
\begin{description}
    \item[img\_processing.py]
        funkcje przetwarzania obrazu i~wczytywania danych,
    \item[blob\_finder.py]
        funkcja wykrywania detali na zdjęciach termowizyjnych,
    \item[blob\_series\_tracker.py]
        funkcje śledzenia i~zliczania detali w~serii zdjęć,
    \item[neural\_network.py]
        funkcje dostępu do modelu sieci i~metod walidacji,
    \item[pliki demo.py]
        pliki demonstrujące działanie utworzonego oprogramowania,
    \item[thesis\_generator.py]
        funkcje generujące obrazy, wykresy oraz tabele,
    \item[docs] dokumentacja projektu.
\end{description}

Projekt zaopatrzono w~plik \texttt{requirements.txt} zawierający listę
wymaganych modułów bibliotecznych.
Stworzone programy najlepiej uruchamiać w~środowisku wirtualnym języka Python
\emph{virtualenv}%
\footnote{%
    Dokumentacja narzędzia virtualenv:
    \url{https://virtualenv.pypa.io/en/latest/}}.
Takie rozwiązanie pozwala na izolację projektu od innych modułów Pythona
zainstalowanych na komputerze oraz daje pewność zgodności wersji używanych
bibliotek.
Programując w~języku Python starano się utrzymać zgodność ze standardem
\emph{PEP8}%
\footnote{%
    Strona standardu PEP8: \url{https://www.python.org/dev/ peps/pep-0008/}}
zapewniającym wysoką jakość i~czytelność kodu.

\section{Dokumentacja projektu}
Za pomocą modułu Pydoc wygenerowano dokumentację projektu w~postaci plików
html.
Zawierają one listę utworzonych funkcji oraz ich opisy.
Pliki dokumentacji zawarto w~repozytorium projektu.

\section{Narzędzia pomocnicze}
W~czasie programowania oraz pisania pracy korzystano z~systemu wersjonowania
Git.
Używanie systemu kontroli wersji zapewnia porządek w~rozwoju projektu oraz
chroni przed przypadkową utratą postępów.
Aby zapewnić odpowiednią jakość oprogramowania użyto programu Pylint%
\footnote{%
    Strona projektu Pylint: \url{https://www.pylint.org}}
(\emph{lintera} języka Python), analizującego kod w~celu wykrycia błędów.
Dokument opisujący projekt inżynierski napisano z~użyciem zestawu makr LaTex
oraz systemu bibliografii BibTex.

\chapter{Spis wykorzystanych funkcji bibliotecznych}
\begin{description}
    \item[Biblioteka standardowa języka Python] moduły: \texttt{glob},
          \texttt{math} oraz \texttt{natsort}.
    \item[Keras] funkcje i~klasy z~przestrzeni nazw \texttt{keras}.
    \item[Matplotlib] funkcje z~przestrzeni nazw \texttt{mpatches} oraz
          \texttt{plt}.
    \item[Numpy] funkcje z~przestrzeni nazw \texttt{np}.
    \item[Pillow] funkcje z~przestrzeni nazw \texttt{Image}.
    \item[Pytesseract] funkcje z przestrzeni nazw \texttt{pytesseract}.
    \item[Scikit-image] funkcje: \texttt{blob\_dog}, \texttt{blob\_doh},
          \texttt{blob\_log}, \texttt{crop}, \texttt{histogram},
          \texttt{img\_as\_ubyte}, \texttt{imread}, \texttt{imsave},
          \texttt{invert}, \texttt{rescale}, \texttt{rgb2gray} oraz
          \texttt{threshold\_otsu}.
    \item[Scikit-learn] funkcja \texttt{train\_test\_split} oraz klasa 
          \texttt{StratifiedKFold}.
    \item[TensorFlow] funkcje i~klasy z~przestrzeni nazw \texttt{tf}.
    \item[Tikzplotlib] funckje z przestrzeni nazw \texttt{tikzplotlib}.   
\end{description}

\end{appendices}

\newpage

% Lists of objects
\listoffigures
\listoftables
\listoflistings

% Bibliography
\nocite{*}
\bibliographystyle{plplain}
\bibliography{references}

\end{document}
